\documentclass{jsarticle}
\usepackage[dvipdfmx,final]{graphicx}
\usepackage{amsmath}
\usepackage{amssymb}
\begin{document}
\title{Zeta function astimate with prime number and number theory \\
differential geometry of quantum level from Morse theory}
\author{Masaaki Yamaguchi}
\date{}
\maketitle

\includegraphics[width=8cm,clip]{D-brane}

   $$ x \to a \to {a \over \log x} \to {x \log x} \to  
    \to \int x \log x dx \to \int{1 \over (x \log x)}dx \to $$
   $$ \to \int \int{1 \over (x \log x)^2}dx_m \to 
    \to {d \over df} \int \int{1 \over (x \log x)^2}dx_m $$
   $$ {d \over df}F = F^{(f)^{'}} = e^f + e^{-f} \to 2i \sin(i x \log x) 
    \to {d \over d\gamma}\Gamma^{-1} = e^f - e^{-f} \to 2 \cos (i x \log x) $$
  　整数から素数の散らばっいる数値をグラフ化する。統合すると、
  ゼータ関数の性質がルールの規則となり、この規則がグラフを変形させたり、
  数式をグラフ化するのに、数式自体を自在に別のこのゼータ関数として成り立つルール
  の範囲で、数式自体を変形して同型と準同型の範囲で数式処理すると、
  TupleSpaceのOmega Databaseにアクセスできてフェルミオンとボソンの
  設計図としてのゼータ関数を求めることができる。
  微分幾何の量子化は、D-braneの単体量からの図形のメッシュ生成として、
  Morse理論となっている。
  

\end{document}
